Defininiamo gli infiniti in $\mathbb{R}$ nel seguente modo.

\begin{definition}
    \textbf{Infinito positivo}

    Sia $\mathbb{A}$ un qualsiasi sottoinsieme di $\mathbb{R}$ non superiormente limitato, allora
    \[ +\infty := \sup\mathbb{A}. \]
\end{definition}

\begin{definition}
    \textbf{Infinito positivo}

    Sia $\mathbb{A}$ un qualsiasi sottoinsieme di $\mathbb{R}$ non inferiormente limitato, allora
    \[ -\infty := \inf\mathbb{A}. \]
\end{definition}

\begin{observation}
    Sia $\mathbb{S}$ un sottoinsieme di $\mathbb{R}$ non superiormente limitato, allora
    l'insieme dei maggioranti di $\mathbb{S}$, $\magg{\mathbb{S}} = \emptyset$.

    Ed essndo $\mathbb{S}\subset \mino{\magg{\mathbb{S}}}$, si può quindi ragionevolmente
    supporre che \[+\infty = \min\emptyset. \]
\end{observation}

\begin{observation}
    Sia $\mathbb{S}$ un sottoinsieme di $\mathbb{R}$ non inferiormente limitato, allora
    l'insieme dei minoranti di $\mathbb{S}$, $\magg{\mathbb{S}} = \emptyset$.

    Ed essndo $\mathbb{S}\subset \magg{\mino{\mathbb{S}}}$, si può quindi ragionevolmente
    supporre che \[-\infty = \max\emptyset.\]
\end{observation}

\begin{proposition}
    $\sup\mathbb{N} = +\infty$ ($\mathbb{N}$ non è superiormente limitato)
\end{proposition}

\begin{demonstration}
    \[ \Th \nexists \Lambda\in\mathbb{R} \mid \Lambda=\sup\mathbb{N} \]

\begin{enumerate}[label=(\roman*)]
    \item Supponiamo per assurdo che 
        \[
            \exists \Lambda\in\mathbb{R} \mid \Lambda=\sup\mathbb{N}
        \]
    \item Essendo $\Lambda$ il minimo dei maggioranti (poiché estremo superiore),
    allora $\Lambda - 1$ non può essere un maggiorante per $\mathbb{N}$, ovvero
        \[
            \exists n\in\mathbb{N}\mid\Lambda-1<n,
            \exists n\in\mathbb{N}\mid\Lambda<n+1.
        \]

    \item Ma sapendo per \defref{def:naturals} che $n\in\mathbb{N}\Rightarrow n+1\in\mathbb{N}$,
    allora
        \[
            \exists n\in\mathbb{N}\mid\Lambda<n
        \]
    \item Dunque $\Lambda=\sup\mathbb{N}$ non può essere un maggiorante per $\mathbb{N}$, non può quindi
    nemmeno essere l'estremo superiore.

    \item La negazione della tesi porta a una contraddizione.
    \item La tesi $\nexists \Lambda\in\mathbb{R} \mid \Lambda=\sup\mathbb{N}$ deve essere
    necessariamente vera.
    \item Quindi $\mathbb{N}$ non è superiormente limitato.
    \item Dunque $\sup\mathbb{N} = +\infty$.
\end{enumerate}
\qed
\end{demonstration}

Definiamo ora le proprietà caratteristiche degli estremi di un qualsiasi
sottoinsieme $\mathbb{S}$ di $\mathbb{R}$.

Sia dunque $\mathbb{S} \subset \mathbb{R}$, allora possiamo introdurre le seguenti definizioni.

\begin{definition}\label{def:sup-lim-rea}
    \textbf{Insieme superiormente limitato}

    Un sottoinsieme $\mathbb{S}$ di $\mathbb{R}$ si dice superiormente limitato con
    $\Lambda=\sup\mathbb{S}$ se e solo se
    \[
        \begin{cases}
            \forall x\in\mathbb{S}, \Lambda \geq x\\
            \forall \epsilon > 0, \exists x\in\mathbb{S} \mid \Lambda-\epsilon<x\leq\Lambda
        \end{cases}
    \]
\end{definition}

\begin{definition}
    \textbf{Insieme inferiormente limitato}

    Un sottoinsieme $\mathbb{S}$ di $\mathbb{R}$ si dice inferiormente limitato con
    $\lambda=\inf\mathbb{S}$ se e solo se
    \[
        \begin{cases}
            \forall x\in\mathbb{S}, \lambda \leq x\\
            \forall \epsilon > 0, \exists x\in\mathbb{S} \mid \lambda\leq x<\lambda+\epsilon
        \end{cases}
    \]
\end{definition}

\begin{example}\leavevmode

    Prendiamo $\mathbb{S} = \left\{ n\in\mathbb{N} \mid 1 - \frac{1}{n} \right\}$

    Considerando un generico $n\in\mathbb{N}$, avremo che
    \[ n > 0, \]
    \[ \frac{1}{n}>0, \]
    \[ -\frac{1}{n} < 0, \]
    \[ 1-\frac{1}{n} < 1, \]
    \[ 1-\frac{1}{n} \leq 1. \]

    Dunque
    \[ \forall n\in\mathbb{N}, 1-\frac{1}{n} \leq 1. \]

    Introducendo anche un gnenerico
    $\epsilon > 0$ e prendendo arbitrariamente $n_\epsilon > \frac{1}{\epsilon}$,
    avremo inoltre che
    \[ n_\epsilon > \frac{1}{\epsilon}, \]
    \[ n_\epsilon \epsilon > 1, \]
    \[ 1 < n_\epsilon \epsilon, \]
    \[ \frac{1}{n_\epsilon} < \epsilon, \]
    \[ 0 \leq \frac{1}{n_\epsilon} < \epsilon, \]
    \[-\epsilon < -\frac{1}{n_\epsilon} \leq 0,\]
    \[1-\epsilon < 1-\frac{1}{n_\epsilon} \leq 1.\]

    Dunque
    \[\forall \epsilon>0 \exists n\in\mathbb{N} \mid 1-\epsilon < 1-\frac{1}{n_\epsilon} \leq 1.\]

    Abbiamo quindi complessivamente che

    \[
        \begin{cases}
            \forall n\in\mathbb{N}, 1 - \frac{1}{n} \leq 1 \\
            \forall \epsilon > 0, \exists n\in\mathbb{N} \mid 1-\epsilon < 1 - \frac{1}{n} \leq 1
        \end{cases}.
    \]

    Ma sapendo che $\forall n\in\mathbb{N}, \exists x=1-\frac{1}{n} | x\in\mathbb{S}$
    e prendendo $\Lambda = 1$, allora

    \[
        \begin{cases}
            \forall x\in\mathbb{S}, x \leq \Lambda \\
            \forall \epsilon > 0, \exists x\in\mathbb{S} \mid 1-\epsilon < x \leq \Lambda
        \end{cases}.
    \]

    Di conseguenza per \defref{def:sup-lim-rea} possiamo dire che
    \[ \sup\mathbb{S} = 1. \] 
\end{example}